\documentclass[book.tex]{subfiles}

\begin{document}

\chapter{Phonon Density Prediction}
\label{chap:phonon_density}
\noindent\rule{11cm}{0.4pt} \\
\textbf{Prerequisites:} \autoref{chap:phonons}, \autoref{chap:molecular_ml}   \\
\textbf{Difficulty Level:} ** \\
\noindent\rule{11cm}{0.4pt}
\vspace{5mm}

Phonons play a critical role in condensed matter science, and are powerful tools to model phenomena such as nanoscale heat transport. The phonon density of states $D(\omega)$ is a core materials property linked to a number of physical properties of materials including specific heat, thermal resistance, potential for superconductivity and more. Existing computational techniques struggle to effectively model this property and are either brittle or expensive.
% (\textcolor{red}{TODO: Cite virtual crystal approximation, density functional perturbation theory})

We can use machine learning methods to direct predict phonon density of states directly from molecular structure \cite{chen2021direct}. Low data availability means that exploiting equivariance becomes crucial for achieving performant behavior. 

\section{Featurization}

Crystal structures are converted into periodic graphs where atoms are nodes and edges connect atoms within a specified distance. Edge $e_{ab}$ is a scalar which represents the distance between atoms $a$ and $b$. Atomic features are one hot encodings of atomic type and mass. The feature vector for hydrogen is given by
\begin{align}
    x_H &= [m_H,0,0,\dotsc] \\
    x_{Li} &= [0, 0, m_{Li},0,\dotsc]
\end{align}

\section{Model Construction}

Equivariant models are constructed to process the periodic graphs and predict phononic density.

Models are trained to predict phonon density of states, a flat vector and are trained on the mean squared loss against DFT data. A diagram of the architecture is shown below

\begin{figure}
    \centering
    % \includegraphics{figures/Differentiable Physics/Materials ML/Phonon Density Prediction/phonon_e3nn.png}
    \includegraphics{figures/Differentiable Physics/Materials ML/Phonon Density Prediction/phonon_den_pred.png}
    \caption{An equivariant architecture for phonon density prediction.}
    \label{fig:enter-label}
\end{figure}
%  https://arxiv.org/pdf/2009.05163.pdf
%"""
%Machine learning has demonstrated great power in materials design,
%discovery, and property prediction. However, despite the success of machine learning in predicting discrete properties, challenges remain for continuous property prediction. The challenge is aggravated in crystalline solids due to crystallographic symmetry considerations and data scarcity. Here, the direct prediction of phonon density-of-states (DOS) is demonstrated using only atomic species and positions as input. Euclidean neural networks are applied, which by construction are equivariant to 3D rotations, translations, and inversion and thereby capture full crystal symmetry, and achieve high-quality prediction using a small training set of $\approx$ 103 examples with over 64 atom types. The predictive model reproduces key features of experimental data and
%even generalizes to materials with unseen elements, and is naturally suited to efficiently predict alloy systems without additional computational cost. The potential of the network is demonstrated by predicting a broad number of high phononic specific heat capacity materials. The work indicates an efficient approach to explore materials’ phonon structure, and can further enable rapid
%screening for high-performance thermal storage materials and
%phonon-mediated superconductors.
%"""

%References 

\section{Materials Design Applications}
A phonon density prediction model can be used to facilitate materials design. The phononic specific heat capacity is given by the formula

\begin{align}
C_V(T) &= \frac{k_B}{m_{tot}} \int_0^\infty \left ( \frac{\hbar \omega}{2 k_B T} \right )^2 \mathrm{csch}^2 \left (\frac{\hbar \omega}{2 k_B T} \right ) g(\omega) d\omega
\end{align}
Here $m_{tot}$ is the mass of all atoms inside the unit cell, and $g(\omega)$ is the normalized phonon density of states. Predicted $g(\omega)$ values can be plugged into this formula to design high specific heat materials.


\end{document}